% ID: SEL-D0073-OB-E01
% Tanda S073.
Sean $x$ e $y$ las cantidades de hornadas A y B. Al expresar el consumo de cada ingrediente y simplificar las dos primeras restricciones, obtenemos
\[
 \begin{aligned}
 x+2y&\le50 &&\text{(harina)},\\
 2x+y&\le44 &&\text{(azúcar)},\\
 x+y&\le26 &&\text{(mantequilla)},\\
 x,y&\ge0.
 \end{aligned}
\]
El beneficio es $B=20x+30y$. La región factible es un polígono cerrado y acotado, por lo que basta evaluar la función lineal en sus vértices. En los ejes se obtienen $(0,0)$, $(22,0)$ y $(0,25)$. La recta de la mantequilla corta a la del azúcar en $(18,8)$ y a la de la harina en $(2,24)$. Las rectas de la harina y del azúcar se cortan en $(38/3,56/3)$, punto que incumple $x+y\le26$.
\[
 \begin{array}{c|rrrrr}
 (x,y)&(0,0)&(22,0)&(18,8)&(2,24)&(0,25)\\
 B\text{ en euros}&0&440&600&760&750
 \end{array}
\]
El máximo se alcanza con $\boxed{2\text{ hornadas A y }24\text{ hornadas B}}$ y da $760$ euros. Como el óptimo ya tiene coordenadas enteras, también es válido para hornadas completas.
